\section{视觉概念设计}
空中机器人的视觉部分承担着整机自主飞行的核心感知与导航职能。与地面哨兵机器人的视觉系统以“敌方装甲板检测—云台跟随—自主射击”为主线不同，空中机器人不依赖对敌方装甲板的检测打击，其视觉系统的首要目标是在比赛场地有限的空域与严格的飞行安全约束下，为飞行控制系统（飞控）提供连续、可靠、低延迟的自身位姿与导航指令，实现“安全飞行，永不坠机”。据此，我们将视觉（导航）系统定义并分解为自定位与状态估计、飞行指令生成、安全保护与故障诊断、数据转发与飞控协同四类核心任务，其功能需求如\tabref{tab:飞机视觉系统功能设计}所示。

\begin{longtable}[htbp!]{ m{5cm}<{\centering}  m{8cm}<{\centering} }
\hline
种类  &  指标 \\
\hline
\endfirsthead
\hline
种类  &  指标 \\
\hline
\endhead
\hline
\multicolumn{2}{r}{接下页}
\endfoot
\caption{飞机视觉（导航）系统功能设计}
\label{tab:飞机视觉系统功能设计}
\endlastfoot
 \hline
 自定位与状态估计  & 输出频率高、延迟低、全局一致的位姿（位置、姿态、速度）；长期运行漂移小，满足空域作业精度 \\  
 \hline 
 飞行指令生成 & 支持航点巡航、方向机动、定点悬停、反反无人机机动等多种飞行模式；指令平滑、满足机体动力学约束，无速度阶跃 \\  
 \hline
 安全保护与故障诊断  & 实时监测定位与指令健康度；定位失效、指令失联、越界时自动悬停/降级/夺权，杜绝失控坠机  \\  
 \hline
 数据转发与飞控协同  & 完成雷达机体坐标系到飞控坐标系（FLU$\rightarrow$ENU/NED）的变换，以飞控可用的形式下发位姿与指令，保证感知-控制闭环 \\  
 \hline
\end{longtable}

\subsection{导航任务定义}
\begin{itemize}
    \item \textbf{自定位与全局状态估计}：融合机载 3D 激光雷达（Livox MID-360）与惯性测量单元（IMU），基于激光惯性里程计算法（FAST-LIO）输出高频、低延迟的机体六自由度位姿与速度。由于激光雷达点云帧率较低（约$\qty{10}{\Hz}$），进一步利用 IMU 预积分对里程计进行插值，将控制所需的位姿反馈提升到$\qty{100}{\Hz}$以上，保证飞行控制的实时性；
    \item \textbf{飞行指令生成与航迹控制}：根据上层决策（电控协议/上位机调试器）给出的目标航点、方向指令或反反无人机指令，结合当前位姿与机体的速度/加速度/角速度约束，生成平滑的速度-位置指令并下发至飞控，实现航点巡航、方向机动、定点悬停与反反无人机垂直振荡等飞行模式；
    \item \textbf{安全保护与故障诊断}：对定位数据的健康度进行实时监测，包括数值有限性校验、坐标越界校验、帧间跳变检测与帧超时检测；一旦判定定位不可信，即停止向飞控转发视觉位姿并触发保护动作。同时通过统一安全管线对输出指令做限幅、加速度约束与地理围栏约束，从算法层杜绝超速、越界与速度阶跃；
    \item \textbf{数据转发与飞控协同}：将视觉定位结果完成从雷达机体坐标系（FLU）到飞控期望坐标系（ENU，飞控内部再转换至 NED）的坐标变换，并以 MAVROS 所要求的话题形式转发至飞控，与飞控的 OFFBOARD 模式闭环协同，实现真正的“视觉-飞控”一体。
\end{itemize}

\subsection{视觉与控制协同概念}
视觉（导航）模块不单独工作，而是与飞控、电控状态机及遥控器权限体系协同，形成“感知—定位—决策—执行—安全”的完整闭环：
\begin{itemize}
    \item \textbf{感知层}：激光雷达与 IMU 原始数据经 FAST-LIO 融合，输出全局里程计与定位状态标志（定位可用/导航可用）；
    \item \textbf{决策与执行层}：导航控制节点（\lstinline{out_control}）根据协议指令或遥控器拨杆选择飞行模式与目标，调用制导与安全管线生成控制指令，经 MAVROS 下发至飞控（PX4）OFFBOARD 模式执行；
    \item \textbf{安全层}：飞行状态监测节点（\lstinline{flight_status}）将飞控连接、位姿可用、导航可用三级状态汇总并上报电控；遥控器摇杆监测节点（\lstinline{rc_stick_monitor}）在 OFFBOARD 模式下检测到遥控器摇杆移动时主动请求切换至 ALTCTL 模式，将控制权交还操作手，实现一键夺权兜底。
\end{itemize}
通过上述多级协同，视觉系统在任何单点失效（定位漂移、指令失联、链路中断）时都能快速降级到安全的悬停或人工接管状态，从根本上保障空中机器人运行的安全性。
